\scanpage{31}

l'aide de la notion de \uline{profondeur}.

Les résultats \inserted{ainsi} obtenus sont l'outil technique principal pour un \deleted{\unsure{étude|dévelop}} \inserted{développement} systématique de \deleted{propriétés} \inserted{certains théorèmes} du type de ``Lefschetz'', comparant les propriétés géométriques d'un schéma $X$ à celles de certains \inserted{sous-schémas} diviseurs $Y$ \deleted{(sections hyperplanes)} \deleted{\unsure{en}} \inserted{sur $X$}, \uline{via} les propriétés du complété formel $X_{/Y}$. \deleted{\unsure{du}} \deleted{Ils \unsure{précisent|prennent} \unsure{une|en} \unsure{remarque|nouvelle} l'observation \unsure{faite|fait} \unsure{par} \unsure{application|appliqué}} Ils sont applicables par exemple \inserted{(\unsure{après des procédés connus|après du procédé général} de passage aux cônes \unsure{projetants|projetés})} au cas d'une variété projective $X$ et une section hyperplane convenable, moyennant des hypothèses \deleted{de reg\unsure{u}} convenables sur la profondeur de $X$. \deleted{\unsure{Ce qui précède} \unsure{Ceci est}}

Dans ce cas \deleted{\unsure{assez voisin du} à préciser} \deleted{la théorie} \deleted{\unsure{le reste}} \deleted{a été} esquissé par H. Grauert qui avait d'ailleurs, \deleted{\unsure{une}} (\unsure{observé|associé|Bourbaki} \deleted{\unsure{à}} \inserted{dès} 1958, (utilisant les th. de \deleted{\unsure{co}} nullité de Serre $[\,]$)$]$ \deleted{de Grauert que dans la situation précédente \unsure{et}} \deleted{\unsure{moyennant des hyp|la situation au-dessus}} \deleted{\unsure{supplémentaires}} \deleted{\unsure{de}}

(i) \inserted{Pour tt $F$ loc. libre,} $H^i(X,F)\to H^i(\hat X,\hat F)$ est bijectif pour des \deleted{petits} $i$ $\leq n$ $[n$ \deleted{dépendant de} \inserted{$\leq\operatorname{prof}F-2$}$]$

(ii) \unsure{d'une|d'un} part \unsure{si} un faisceau formel loc libre $\mathcal F$ sur $X_{/Y}$, \deleted{$\hat{\mathcal F}$ \unsure{est} \unsure{isom}} \inserted{et si $\operatorname{prof}(\mathcal O_{\hat X})\geq 2$,} alors pour \unsure{$n$|un} \deleted{coupe de sections ``} \unsure{grand|formel}, $\hat{\mathcal F}\unsure{(n)}$ est \unsure{engendré|engendre} par ses sections [$\unsure{\text{sauf}|\exists n}$] sur un ens fini -- $[$et \deleted{\unsure{à preuve le verrez|ce, sous les autres}} \inserted{en fait} \unsure{pratiquement|pendant} \deleted{\unsure{je}} comme \unsure{nous montrerons|une immersion}$]$, \deleted{Alors} \deleted{\unsure{des conditions|des Grothendieck}} \deleted{\unsure{C'est cette une des}} (Cela qui impliquait \inserted{p.ex.} déjà le th de Lefschetz pour le groupe fondamental énoncé en $[\,]$.)\margin{La théorie développée ici a pour but de préciser et de \deleted{\unsure{f}} généraliser les résultats de Grauert.}
% margin note written sideways

% archive sheet number: 30
